\begin{table}[htbp]
\centering
\caption{Table A1: Summary Statistics (District Level)}
\label{tab:fetzer_a1}
\begin{adjustbox}{max width=\textwidth}
\begin{tabular}{lcccc}
\toprule
Variable & Obs & Mean & Std.\ Dev. & Min & Max \\
\midrule
`var' ---"
  3.     sum `var'
  4. }

--- Summary: pct_votes_UKIP & 3290 & 4.454 & 7.571 & 0.000 & 44.261 \\
EP \% for UKIP & 1140 & 21.118 & 9.397 & 3.596 & 51.581 \\
`var' (year==2011) ---"
  3.     sum `var' if year==2011
  4. }

--- Summary: QUAL_ALL_noq_sh & 346 & 0.286 & 0.062 & 0.130 & 0.456 \\
\% Routine occupations (2001) & 346 & 0.102 & 0.030 & 0.036 & 0.206 \\
\% in Retail (2001) & 346 & 0.169 & 0.021 & 0.100 & 0.241 \\
\% in Manufacturing (2001) & 346 & 0.154 & 0.054 & 0.054 & 0.337 \\
Total austerity impact & 378 & 447.122 & 121.110 & 177.000 & 914.000 \\
Tax credits impact & 379 & 87.971 & 23.563 & 0.000 & 150.000 \\
Child benefit impact & 379 & 71.517 & 9.425 & 37.000 & 102.000 \\
Council tax benefit impact & 379 & 7.211 & 7.810 & 0.000 & 40.000 \\
Disability living allowance impact & 379 & 36.570 & 12.204 & 14.000 & 80.000 \\
Bedroom tax impact & 379 & 10.813 & 5.597 & 2.000 & 30.000 \\
\bottomrule
\end{tabular}
\end{adjustbox}
\parbox{\textwidth}{\footnotesize\textit{Notes:} District-level summary statistics. Austerity variables measured as financial loss per working-age person per year (\pounds).}
\end{table}